\documentclass[../../main.tex]{subfiles}% 注意这里的文件路径不能用 ./main.tex ,否则用latexmk编译子文件会报错
\graphicspath{{\subfix{./image/}}{\subfix{../../image/}}} % 指定图片引用路径,后续可以直接使用图片文件名
% 如果图片存放在上级目录,则使用 ../../image/ 作为备用路径

% 例如:
% \begin{figure}[H]
% \centering
% \includegraphics[width=0.3\textwidth]{图.png}
% \caption{}
% \label{figure:图}
% \end{figure}
% 注意:上述\label{}一定要放在\caption{}之后,否则引用图片序号会只会显示??.

\begin{document}

\section{Cantor(三分)集}

\begin{definition}[Cantor三分集]\label{definition:Cantor三分集}
设 \( I_0 = \{[0,1]\} \)。对 \( n \geqslant 0 \)，假定 \( I_n \) 是由 \( 2^n \) 个互不相交的闭区间组成的集合。定义
\[
d_n = \frac{1}{3^{n+1}}.
\]
对每个区间 \([a,b] \in I_n\)，移除其中点附近长度为 \( d_n \) 的开区间
\[
\left( \frac{a+b}{2} - \frac{d_n}{2},\; \frac{a+b}{2} + \frac{d_n}{2} \right),
\]
剩余两个闭区间
\[
\left[ a,\; \frac{a+b}{2} - \frac{d_n}{2} \right] \quad\text{和}\quad \left[ \frac{a+b}{2} + \frac{d_n}{2},\; b \right],
\]
记所有这样得到的新区间构成的集合为 \( I_{n+1} \)。由此可得到一列集合 \(\{I_n\}_{n=0}^{\infty}\)。
令
\[
C_n = \bigcup_{I \in I_n} I, \qquad n = 0,1,2,\cdots
\]
则 \textbf{Cantor 三分集} 定义为
\[
\mathcal{C} = \bigcap_{n=0}^{\infty} C_n.
\]
\end{definition}
\begin{remark}
序列 \(\{I_n\}_{n=0}^{\infty}\) 的良定义性证明：令
\[
S = \bigcup_{k=0}^{\infty} \left\{ I \subseteq \mathcal{P}([0,1]) \;\middle|\; I \text{ 是由 } 2^k \text{ 个互不相交的闭区间组成的集合} \right\}.
\]
取 \( I_0 = \{[0,1]\} \in S \)。对\(\forall n \in \mathbb{N},\forall I \in S \)，存在唯一 \( k \in \mathbb{N} \) 以及序列 \(\{a_i\}_{i=1}^{2^k},\{b_i\}_{i=1}^{2^k}\) 使得
\[
a_i < b_i,\quad b_i < a_{i+1},\quad I = \bigcup_{i=1}^{2^k} \{\,[a_i,b_i]\,\}.
\]
记 \( d_n = \frac{1}{3^{n+1}} \)。对每个区间 \([a_i, b_i]\)移除开区间
\[
\left( \frac{a_i+b_i}{2}-\frac{d_n}{2},\frac{a_i+b_i}{2}+\frac{d_i}{2} \right) ,
\]
剩余两个闭区间
\[
\left[ a_i,\;\frac{a_i+b_i}{2}-\frac{d_n}{2} \right] \quad \text{和}\quad \left[ \frac{a_i+b_i}{2}+\frac{d_n}{2},\;b_i \right] .
\]
令
\[
I_n' =\bigcup_{i=1}^{2^k}{\left\{ \left[ a_i,\;\frac{a_i+b_i}{2}-\frac{d_n}{2} \right] ,\;\left[ \frac{a_i+b_i}{2}+\frac{d_n}{2},\;b_i \right] \right\}}.
\]
则 \( I_n' \) 是由 \( 2^{k+1} \) 个互不相交的闭区间组成的集合，故 \( I_n' \in S \)。定义 \( f_n(I) = I_n' \)，则
\[
f_n : S \to S,\quad I \mapsto f_n(I) = I_n'
\]
是良定义的。根据\hyperref[Set Theory-theorem:递归定理]{递归定理}，存在唯一的函数 \( \varphi : \mathbb{N} \to S \) 满足
\[
\varphi(0) = I_0,\quad \varphi(n+1) = f_n(\varphi(n))(\forall n \in \mathbb{N}).
\]
对任意 \( n \in \mathbb{N} \)，定义 \( I_n \triangleq \varphi(n) \)。因此，序列 \(\{I_n\}_{n=0}^{\infty}\) 是良定义的。
\end{remark}

\begin{theorem}[Cantor集的基本性质]\label{theorem:Cantor集的基本性质}
设集合$C$是Cantor集,则
\begin{enumerate}[(1)]
\item $C$是非空有界闭集,因此是紧集.

\item $C = C'$,即$C$为完全集.

\item $C$无内点.即$C$是完备集.

\item Cantor集的基数是$c$.

\item $[0,1]\setminus C$的长度的总和为1.

\item 对$\forall n\geqslant 2$,有
\begin{align*}
C_n = [0,1] \setminus \bigcup\limits_{m=1}^n \bigcup\limits_{i=0}^{3^{m-1}-1} \left( \frac{3i+1}{3^m}, \frac{3i+2}{3^m} \right) = \bigcap\limits_{m=1}^n \bigcup\limits_{i=0}^{3^{m-1}-1} \left( \left[ \frac{3i}{3^m}, \frac{3i+1}{3^m} \right] \cup \left[ \frac{3i+2}{3^m}, \frac{3i+3}{3^m} \right] \right),
\end{align*}
其中$C_n$的定义同\refdef{definition:Cantor三分集}.进而
\begin{align*}
[0,1]\setminus C=\bigcup_{m=1}^{\infty}{\bigcup_{i=0}^{3^{m-1}-1}{\left( \frac{3i+1}{3^m},\frac{3i+2}{3^m} \right)}}.
\end{align*}
\end{enumerate}
\end{theorem}
\begin{proof}
\begin{enumerate}[(1)]
\item 因为每个$F_n$都是非空有界闭集，而且$F_n\supset F_{n + 1}$，所以根据Cantor闭集套定理，可知$C$不是空集（实际上，$F_n$ ($n = 1,2,\cdots$) 中每个闭区间的端点都是没有被移去的，即都是$C$中的点）. 显然，$C$是闭集.

\item 设$x\in C$，则$x\in F_n$ ($n = 1,2,\cdots$)，即对每个$n$，$x$属于长度为$1/3^n$的$2^n$个闭区间中的一个. 于是，对任一$\delta>0$，存在$n$，满足$1/3^n<\delta$，使得$F_n$中包含$x$的闭区间含于$(x - \delta, x + \delta)$. 此闭区间有两个端点，它们是$C$中的点且总有一个不是$x$. 这就说明$x$是$C$的极限点，故得$C'\supset C$. 由(i)知$C = C'$.

\item 设$x\in C$，给定任一区间$(x - \delta, x + \delta)$，取$2/3^n<\delta$. 因为$x\in F_n$，所以$F_n$中必有某个长度为$1/3^n$的闭区间$F_{n,k}$含于$(x - \delta, x + \delta)$. 然而，在构造$C$集的第$n + 1$步时，将移去$F_{n,k}$的中央三分开区间. 这说明$(x - \delta, x + \delta)$不含于$C$.

\item 事实上，将$[0,1]$中的实数按三进位小数展开，则Cantor集中点$x$与下述三进位小数集的元
\begin{align*}
x=\sum_{i = 1}^{\infty}\frac{a_i}{3^i}, \quad a_i = 0,2
\end{align*}
一一对应. 从而知$C$为连续基数集（与$(0,1]$的二进位小数比较）. 

\item 由$C$的定义可得
\begin{align*}
\sum_{n = 1}^{\infty}2^{n - 1}3^{-n}=1.
\end{align*}
\end{enumerate}
\end{proof}

\begin{definition}[类Cantor集]
设$\delta$是$(0,1)$内任意给定的数,考虑在$[0,1]$区间, 取$p=\frac{1+2\delta}{\delta}$，采用类似于Cantor集的构造过程:

第一步，移去长度为$\frac{1}{p}$的同心开区间；

第二步，在留存的两个闭区间的每一个中，又移去长度为$\frac{1}{p^2}$的同心开区间；

第三步，在留存的四个闭区间中再移去长度为$\frac{1}{p^3}$的同心区间. 继续此过程，可得一列移去的开区间，记其并集为$G$（开集）,
我们称$C_p=[0,1]\setminus G$为\textbf{类Cantor集}（当$p = 3$时，$C_p$就是Cantor（三分）集）.类Cantor集也称为\textbf{Harnack集}. 
\end{definition}
\begin{remark}
若要在$\mathbb{R}^n$的单位方体$[0,1]\times[0,1]\times\cdots\times[0,1]$中构造具有类似性质的集合，则只需取$C\times C\times\cdots\times C$（$C$是$[0,1]$中的类Cantor集）即可.
\end{remark}

\begin{theorem}[类Cantor集的基本性质]\label{theorem:类Cantor集的基本性质}
\begin{enumerate}[(1)]
\item $G$的总长度为$\delta$（$0<\delta<1$是任意给定的数）的稠密开集.

\item $C_p$是非空完全集，且没有内点.
\end{enumerate}
\end{theorem}
\begin{proof}
\begin{enumerate}[(1)]
\item 由$G$的定义可知$G$的总长度为
\begin{align*}
\sum_{n = 1}^{\infty}2^{n - 1}\left(\frac{1}{p}\right)^n=\frac{1}{p - 2}=\delta.
\end{align*}

\item 
\end{enumerate}
\end{proof}

\begin{definition}[Cantor函数]\label{definition:Cantor函数}
设 \(C\) 是 \([0,1]\) 中的 Cantor 集, 其中的点我们用三进位小数
\begin{align*}
x = 2\sum_{i = 1}^{\infty}\frac{\alpha_i}{3^i}, \quad \alpha_i = 0,1 \quad (i = 1,2,\cdots)
\end{align*}
来表示.

(i) 作定义在 \(C\) 上的函数 \(\varphi(x)\). 对于 \(x \in C\), 定义
\[
\varphi(x)=\varphi\left(2\sum_{i = 1}^{\infty}\frac{\alpha_i}{3^i}\right)=\sum_{i = 1}^{\infty}\frac{\alpha_i}{2^i}, \quad \alpha_i = 0,1 \quad (i = 1,2,\cdots).
\]

(ii) 作定义在 \([0,1]\) 上的 \(\varPhi(x)\). 对于 \(x \in [0,1]\), 定义
\[
\varPhi(x)=\sup\{\varphi(y):y\in C,y\leqslant x\}.
\]
我们称 \(\varPhi(x)\) 为 \textbf{Cantor函数}.
\end{definition}

\begin{theorem}[Cantor函数的性质]\label{theorem:Cantor函数的性质}
设$\varPhi(x)=\sup\{\varphi(y):y\in C,y\leqslant x\}$为Cantor函数, 则有下列性质:
\begin{enumerate}[(1)]
\item \(\varphi(C)=[0,1]\),即$\varphi$是满射.并且\(\varphi(x)\) 是 \(C\) 上的递增函数.

\item \(\varPhi(x)\) 是 \([0,1]\) 上的递增连续函数.此外, 在构造 Cantor 集的过程中所移去的每个中央三分开区间 \(I_{n,k}\) 上, \(\varPhi(x)\) 都是常数. 
\end{enumerate}
\end{theorem}
\begin{proof}
\begin{enumerate}[(1)]
\item 因为 \([0,1]\) 中的点可用二进位小数表示, 所以由$\varphi$的定义有 \(\varphi(C)=[0,1]\)..

下面证明$\varphi(x)$是$C$上的递增函数.设 \(\alpha_1,\alpha_2,\cdots,\beta_1,\beta_2,\cdots\) 是取 \(0\) 或 \(1\) 的数, 而且它们所表示的 \(C\) 中的数有下述关系:
\[
2\sum_{i = 1}^{\infty}\frac{\alpha_i}{3^i}<2\sum_{i = 1}^{\infty}\frac{\beta_i}{3^i}.
\]
若记 \(k = \min\{i:\alpha_i\neq\beta_i\}\), 则我们有
\begin{align*}
0&<\sum_{i = 1}^{\infty}\frac{\beta_i - \alpha_i}{3^i}=\frac{\beta_k - \alpha_k}{3^k}+\sum_{i > k}\frac{\beta_i - \alpha_i}{3^i}\\
&\leqslant\frac{\beta_k - \alpha_k}{3^k}+\sum_{i > k}\frac{2}{3^i}=\frac{\beta_k - \alpha_k + 1}{3^k}.
\end{align*}
由此可知 \((\alpha_k<\beta_k)\alpha_k = 0,\beta_k = 1\), 从而得到
\begin{align*}
\varphi\left(2\sum_{i = 1}^{\infty}\frac{\alpha_i}{3^i}\right)&=\sum_{i = 1}^{\infty}\frac{\alpha_i}{2^i}=\sum_{i = 1}^{k - 1}\frac{\alpha_i}{2^i}+\sum_{i = k}^{\infty}\frac{\alpha_i}{2^i}\\
&\leqslant\sum_{i = 1}^{k - 1}\frac{\beta_i}{2^i}+\sum_{i = k + 1}^{\infty}\frac{1}{2^i}=\sum_{i = 1}^{k - 1}\frac{\beta_i}{2^i}+\frac{1}{2^k}\\
&\leqslant\sum_{i = 1}^{k - 1}\frac{\beta_i}{2^i}+\sum_{i = k}^{\infty}\frac{\beta_i}{2^i}=\varphi\left(2\sum_{i = 1}^{\infty}\frac{\beta_i}{3^i}\right).
\end{align*}

\item 由(1)的结论及$\varPhi$的定义即得$\varPhi$的递增性.因为 \(\varPhi([0,1]) = [0,1]\), 所以由\refpro{Basis of Analytics-proposition:定义是区间的单调函数值域还是区间就必是连续函数}可知\(\varPhi(x)\) 是 \([0,1]\) 上的连续函数. 
\end{enumerate}
\end{proof}

\begin{definition}[Smith-Volterra-Cantor集(SVC集)]\label{definition:Smith-Volterra-Cantor集(SVC集)}
设 \( I_0 = \{[0,1]\} \)。对 \( n \geqslant 0 \)，假定 \( I_n \)是由 \( 2^n \) 个互不相交的闭区间组成的集合,定义
\[
d_n = \frac{1}{2^{2n+2}}.
\]
对每个区间 \([a,b] \in I_n\)，移除其中点附近长度为 \( d_n \) 的开区间
\[
\left( \frac{a+b}{2} - \frac{d_n}{2},\; \frac{a+b}{2} + \frac{d_n}{2} \right),
\]
剩余两个闭区间
\[
\left[ a,\; \frac{a+b}{2} - \frac{d_n}{2} \right]\text{和}\left[ \frac{a+b}{2} + \frac{d_n}{2},\; b \right],
\]
记所有这样得到的新区间构成的集合为 \( I_{n+1} \)。由此可得到一列集合$\{I_n\}_{n=0}^{\infty}$.
令
\[
S_n \triangleq \bigcup_{I \in I_n} I, \qquad n = 0,1,2,\cdots
\]
则 \textbf{Smith–Volterra–Cantor 集}（简称 $\mathbf{SVC}$\textbf{集}）定义为
\[
S \triangleq \bigcap_{n=0}^{\infty} S_n  .
\]
\end{definition}
\begin{remark}
序列\(\{I_n\}_{n=0}^{\infty}\)的良定义性证明:令
\[
S = \bigcup_{k=0}^{\infty} \left\{ I \subseteq \mathcal{P}([0,1])  \middle| I \text{是由} 2^k \text{ 个互不相交的闭区间组成的集合} \right\}.
\]
取$I_0=\{[0,1]\}$.对\(\forall n \in \mathbb{N}\)，\(\forall I \in S\),存在唯一 \(k\in \mathbb{N}\)和$\left\{ a_i \right\} _{i=1}^{2^k},\left\{ b_i \right\} _{i=1}^{2^k}$,使得
\[
a_i<b_i,\quad b_i<a_{i+1},\quad I=\bigcup_{i=1}^{2^k}{\{\left[ a_i,b_i \right]\}}.
\]
记$d_n = \frac{1}{2^{2n+2}}$,
对每个区间 \([a_i, b_i]\)移除开区间
\[
\left( \frac{a_i+b_i}{2}-\frac{d_n}{2},\frac{a_i+b_i}{2}+\frac{d_i}{2} \right) ,
\]
剩余两个闭区间
\[
\left[ a_i,\;\frac{a_i+b_i}{2}-\frac{d_n}{2} \right] \quad \text{和}\quad \left[ \frac{a_i+b_i}{2}+\frac{d_n}{2},\;b_i \right] .
\]
令
\[
I_n' =\bigcup_{i=1}^{2^k}{\left\{ \left[ a_i,\;\frac{a_i+b_i}{2}-\frac{d_n}{2} \right] ,\;\left[ \frac{a_i+b_i}{2}+\frac{d_n}{2},\;b_i \right] \right\}}.
\]
则 \(I_n'\) 是由 \(2^{k+1}\) 个互不相交的闭区间构成的集合，故 \(I_n' \in S\)。定义$f_n(I) = I_n'$,则
\begin{align*}
f_n:S\to S,\,\,I\mapsto f_n(I)=I_n'
\end{align*}
是良定义的.
根据\hyperref[Set Theory-theorem:逆归定理]{递归定理}，存在唯一的函数 \(\varphi : \mathbb{N} \to S\) 满足
\[
\varphi(0) = I_0,\quad \varphi(n+1) = f_n(\varphi(n))(\forall n \in \mathbb{N}).
\]
对\(\forall n \in \mathbb{N}\)，定义$I_n \triangleq \varphi(n).$
因此，序列 \(\{I_n\}_{n=0}^{\infty}\) 是良定义的.
\end{remark}

\begin{theorem}\label{theorem:SVC集的基本性质}
设集合$S$是SVC集,则
\begin{enumerate}[(1)]
\item $S$是闭的.

\item $S$是无处稠密的.

\item $S$的Lebesgue测度为$m(S)=\frac{1}{2}$.

\item\label{theorem:SVC集的基本性质-4} 定义$I_n,d_n(\forall n\in \mathbb{N})$同\refdef{definition:Smith-Volterra-Cantor集(SVC集)},则
\begin{align*}
\left[ 0,1 \right] \backslash S=\bigcup_{n=0}^{\infty}{\bigcup_{\left[ a,b \right] \in I_n}{\left( \frac{a+b}{2}-\frac{d_n}{2},\;\frac{a+b}{2}+\frac{d_n}{2} \right)}}.
\end{align*}
\end{enumerate}
\end{theorem}
\begin{proof}
\begin{enumerate}[(1)]
\item 

\item 

\item 由SVC集的定义有
\[
m(S) = 1 - \sum_{n=0}^{\infty} 2^n \frac{1}{2^{2n+2}} = 1 - \sum_{n=0}^{\infty} \frac{1}{2^{n+2}} = \frac{1}{2}.
\]
\end{enumerate}
\end{proof}

\begin{definition}[Volterra(沃尔泰拉)函数]\label{definition:Volterra(沃尔泰拉)函数}
设集合$S$是SVC集,令 \( \{ (a_k, b_k) \}_{k=1}^{\infty} \) 为构造 \( S \) 时移除的所有开区间,由\rrefthe{theorem:SVC集的基本性质}{theorem:SVC集的基本性质-4}知,此即
\[ 
\bigcup_{n=0}^{\infty} \bigcup_{[a,b] \in I_n} \bigl( \frac{a+b}{2} - \frac{d_n}{2},\; \frac{a+b}{2} + \frac{d_n}{2} \bigr) 
\]中的所有开区间,其中$I_n,d_n(\forall n\in \mathbb{N})$的定义同\refdef{definition:Smith-Volterra-Cantor集(SVC集)}.定义函数 \( V : [0,1] \to \mathbb{R} \) 如下：
\[
V(x) =
\begin{cases}
0, & x \in S, \\[6pt]
(x - a_k)^2 \sin\!\left(\dfrac{1}{x - a_k}\right), & x \in \left(a_k,\; \dfrac{a_k+b_k}{2}\right], \\[10pt]
(b_k - x)^2 \sin\!\left(\dfrac{1}{b_k - x}\right), & x \in \left[\dfrac{a_k+b_k}{2},\; b_k\right).
\end{cases}
\]
函数 \( V \) 称为 \textbf{Volterra 函数}。
\end{definition}

\begin{theorem}[Volterra函数的基本性质]\label{theorem:Volterra函数的基本性质}
设$V$是Volterra函数,则
\begin{enumerate}[(1)]
\item \( V \) 在 \( [0,1] \) 上处处可导；
\item \( V \)的导函数 \( V' \) 在 \( S \) 上的每一点都不连续，而 \( S \) 具有正测度，因此 \( V' \) 在 \( [0,1] \) 上不是 Riemann 可积的。
\end{enumerate}
进而$V'$在$[0,1]$上存在原函数$V$,但$V'$在$[0,1]$上不是 Riemann 可积的.
\end{theorem}

\begin{example}
\(E\subset\mathbb{R}\) 是完全集当且仅当 \(E = \left(\bigcup_{n\geqslant 1}(a_n,b_n)\right)^c\)，其中 \((a_i,b_i)\) 与 \((a_j,b_j)\) \((i\neq j)\) 无公共端点。
\end{example}
\begin{proof}

{\heiti 必要性：}若 \(E\) 是完全集，则 \(E\) 是闭集。从而 \(E^c\) 是开集，它是 \(E^c\) 内构成区间的并集。这些构成区间相互之间是没有公共端点的，否则 \(E\) 中就会有孤立点了，这是不可能的。

{\heiti 充分性：}首先，由题设知 \(E\) 是闭集。其次，对任意的 \(x\in E\)，如果 \(x\notin E'\)，那么存在 \(\delta>0\)，使得 \((x - \delta,x + \delta)\cap E = \{x\}\)。这说明 \(x\) 是某两个开区间的端点，与假设矛盾。
\end{proof}

\begin{example}
设 \(E\subset\mathbb{R}^2\) 是完全集，则 \(E\) 是不可数集。
\end{example}
\begin{proof}
用反证法。假定 \(E = \{x_n\in\mathbb{R}^2: n = 1,2,\cdots\}\)。

(i) 选取 \(y_1\in E\setminus\{x_1\}\)，则点 \(x_1\) 到 \(y_1\) 的距离大于 \(0\)。存在以 \(y_1\) 为中心的闭正方形 \(Q_1\)，\(Q_1\cap E\) 是紧集。

(ii) 看 \(E\setminus\{x_2\}\)。因为 \(y_1\) 是 \(E\setminus\{x_2\}\) 的极限点，所以 \(\mathring{Q}_1\cap (E\setminus\{x_2\})\neq\varnothing\)。又取 \(y_2\in\mathring{Q}_1\cap (E\setminus\{x_2\})\)，并作以 \(y_2\) 为中心的闭正方形 \(Q_2\)：\(Q_2\subset Q_1\)，\(x_1\notin Q_2\)，\(x_2\notin Q_2\)，可知 \((Q_1\cap E)\supset (Q_2\cap E)\) 是紧集。如此继续做下去，可得有界闭集套列 \(\{Q_n\cap E\}\)：\((Q_{n - 1}\cap E)\supset (Q_n\cap E)\) \((n\in\mathbb{N})\)，而且 \(x_1,x_2,\cdots,x_n\) 不在其内。我们有
\begin{align*}
\bigcap_{n = 1}^{\infty}(Q_n\cap E)=\varnothing,
\end{align*}
导致矛盾。 
\end{proof}

\begin{proposition}
任一非空完全集的基数均为 \(c\)。
\end{proposition}
\begin{proof}
证明见那汤松著《实变函数论》的上册，有高等教育出版社出版的中译本，1955 年.
\end{proof}

\begin{example}
设 \(E = \left\{x\in[0,1]: x = \sum_{n = 1}^{\infty}a_n/10^n, a_n = 2 \text{ 或 } 7\right\}\)，我们有

(i) \(E\) 是闭集； 

(ii) \(\overline{E}=c\)； 

(iii) \(E\) 在 \([0,1]\) 中不稠密。
\end{example}
\begin{proof}
(i) 若有 \(\{x_m\}\subset E\)：\(x_m\rightarrow x (m\rightarrow\infty)\)，则
\begin{align*}
x = \sum_{n = 1}^{\infty}b_n/10^n \quad (b_n = 0,1,2,\cdots,9).
\end{align*}
如果 \(|x_m - x|<10^{-p}\)，那么在 \(x\in E\) 时，\(b_n = 2 \text{ 或 } 7 (n = 1,2,\cdots,p - 1)\)。这说明 \(E\) 是闭集。

(ii) 与 \(0\) 和 \(1\) 组成的数列类似，\(\overline{E}=c\)。

(iii) 注意到 \(E\cap(0.28,0.7)=\varnothing\)，故 \(E\) 不是稠密集。 
\end{proof}





\end{document}